\chapter{文献综述}


\section{研究综述}

以往针对数学师范生教学技能的研究丰富，涉及范围广泛，关注了数学教学技能的多个方面。

李渺等把课堂教学技能分为教学语言、板书、讲解、提问、导入和结束六种，并进行调查分析，研究结果认为知识因素和心理因素是影响师范生课堂教学技能提高的主要因素\cite{李渺2009高师数学专业师范生课堂教学技能的调查研究}。

何念如指出了师范生在教案设计、教学方法等方面的不足，并提出了提高数学师范生教学技能的对策\cite{何念如数学师范生教学技能存在的问题及对策研究}。

闫德明等总结出提高数学师范生教学技能的途径的四个特点：所有学生都认为实习和学习最重要、优秀生更重视多种途径、普通生较为重视家教、所有学生都不重视微格教学\cite{闫德明提高数学师范生教学技能途径的调查分析}。

于海燕等提出了促进数学师范生教学技能提升的四个方法：优化课程设置、引进视频案例教学法、开展微格教学、加强教育实践\cite{于海燕2018数学师范生教学技能培养的研究}。

徐国宁提出数学教师应具备以下教学技能：组织学生合作和探究学习的技能；利用和开发课程资源的技能；使用现代教育技术手段及设计微课程和制作微课视频的技能；创新的教学技能；对学生进行科学评价的技能；进行教育科研的技能\cite{徐国宁2019新课程理念下数学教师应具备的教学技能}。



\section{述评}

当前研究涉及到数学师范生教学技能的分类、提高等方面，对个别教学技能研究较多，而对其他教学技能研究较少，并且未进行全面的调查分析，未得出师范生其它数学教学技能发展的不足之处，并且没有深入探究数学教学技能各个维度之间的内在关联。